\section{Проектирование и реализация сервиса}
\thispagestyle{plain}

% TODO: ~16 страниц --- инженерная глава.

\subsection{Архитектура клиент-сервер}
% C4-диаграммы (контекст, контейнеры) обновлённые относительно НИР.
% Принципиальное отличие: серверный компонент собственной разработки
% вместо Cloud Firestore. Обоснование архитектурных решений.

\subsection{Технологический стек}
% Frontend: Flutter 3.x + Riverpod 2.6 + Material Design 3,
% единая кодовая база на Web/iOS/Android/macOS.
% Backend: Go 1.22 + Gin, gorilla/websocket для real-time.
% БД: PostgreSQL 16. Кэш и pub/sub: Redis 7.
% Reverse proxy: Caddy с автоматическим TLS.
% Контейнеризация: Docker + Docker Compose.
% Мониторинг: Prometheus + Grafana.
% Хостинг: Selectel/Timeweb (Российская Федерация).

\subsection{Структура базы данных}
% ER-диаграмма основных коллекций: users, courses, question\_banks,
% questions, rooms, sessions, participants, answers, gamification.
% Схема PostgreSQL с миграциями (golang-migrate).
% Индексы и стратегии партиционирования.

\subsection{REST API и WebSocket-протокол}
% Принципы построения REST API (OpenAPI спецификация).
% WebSocket-протокол для интерактивных сессий: типы событий,
% форматы сообщений, обработка переподключений.
% Аутентификация: JWT-токены.

\subsection{Интеграция с университетской LMS edu.gubkin}
% Архитектура интеграции: импорт банка вопросов в форматах
% Moodle XML / GIFT / QTI в зависимости от поддержки в edu.gubkin.
% Маппинг типов вопросов между LMS и сервисом.
% Сценарий преподавателя: выгрузка из edu.gubkin
% $\rightarrow$ загрузка в EduQuiz $\rightarrow$ настройка геймификации
% $\rightarrow$ проведение сессии.

\subsection{Реализация клиента и серверной части}
% Структура Flutter-проекта: feature-first архитектура.
% Структура Go-проекта: слои handlers/services/repository.
% Ключевые алгоритмы: подсчёт XP, SM-2, лидерборд.
% Адаптивная вёрстка: phone, tablet, desktop, web.

\subsection{Развёртывание}
% Docker Compose для production: api, postgres, redis, prometheus, grafana, caddy.
% Конфигурация Caddy и автоматический выпуск TLS.
% CI/CD на GitHub Actions: сборка, тесты, публикация Docker-образов.
% Развёртывание на VPS Selectel/Timeweb: 2 vCPU, 4 ГБ RAM, 40 ГБ SSD.

\textbf{Выводы по главе 3.}
% Что реализовано, какие технологические решения приняты и почему.

\newpage
